% Luc Destombe --- licence CC BY-NC-SA 4.0
% https://creativecommons.org/licenses/by-nc-sa/4.0/deed.fr
% Fichier autonome : compiler avec pdflatex, deux fois (signets, nombre de pages).
\def\lyceeversion{1.3}
\documentclass[a4paper,10pt]{article}
\makeatletter
\newif\iflycee@niveaux
\newif\iflycee@cm
\newif\iflycee@palatino

\RequirePackage[utf8]{inputenc}
\RequirePackage[T1]{fontenc}
\RequirePackage[french]{babel}
\RequirePackage{etoolbox}

\newcommand{\ShorthandsOff}{\@ifundefined{shorthandoff}{}{\shorthandoff{;:!?}}}
\newcommand{\ShorthandsOn}{\@ifundefined{shorthandon}{}{\shorthandon{;:!?}}}
\ShorthandsOff
\AtBeginDocument{\ShorthandsOn}
\AtBeginEnvironment{tikzpicture}{\ShorthandsOff}

\RequirePackage{geometry}
\RequirePackage{fancyhdr}
\RequirePackage{lastpage}
\RequirePackage{multicol}
\RequirePackage{array}
\RequirePackage{enumitem}
\RequirePackage{xcolor}
\RequirePackage{needspace}

\RequirePackage{amsmath,amssymb}
\iflycee@cm\else\RequirePackage{fourier}\fi
\DeclareMathAlphabet{\mathcal}{OMS}{cmsy}{m}{n}
\renewcommand{\ttdefault}{lmtt}
\begingroup\catcode`\_=\active \catcode`\^=\active
\global\let_\sb \global\let^\sp
\endgroup
\AtBeginDocument{\catcode`\_=12 \mathcode`\_="8000
  \catcode`\^=12 \mathcode`\^="8000 }
\newlength{\retraitcalculs}
\setlength{\retraitcalculs}{2em}

\RequirePackage{tikz}
\usetikzlibrary{arrows.meta,calc,babel,shapes.misc}
\RequirePackage[most]{tcolorbox}
\tcbuselibrary{skins,breakable}

\definecolor{coulDef}{HTML}{1F4E79}
\definecolor{coulProp}{HTML}{7030A0}
\definecolor{coulRem}{HTML}{C55A11}
\definecolor{coulEx}{HTML}{2E7D32}
\definecolor{coulExo}{HTML}{B03A2E}
\definecolor{coulPart}{HTML}{1F4E79}
\definecolor{coulMeth}{HTML}{1B6E4F}
\definecolor{coulCourbe}{HTML}{0B5ED7}
\definecolor{coulCourbeB}{HTML}{C00000}
\definecolor{coulCourbeC}{HTML}{2E7D32}
\definecolor{coulGrille}{HTML}{8C8C8C}
\definecolor{coulRep}{HTML}{1B6E4F}
\definecolor{coulBar}{HTML}{2E7D32}

\newcommand{\R}{\mathbb{R}}
\newcommand{\Rs}{\mathbb{R}^{*}}
\newcommand{\N}{\mathbb{N}}
\newcommand{\Ns}{\mathbb{N}^{*}}
\newcommand{\Z}{\mathbb{Z}}
\newcommand{\Q}{\mathbb{Q}}
\renewcommand{\le}{\leqslant}
\renewcommand{\ge}{\geqslant}
\renewcommand{\approx}{\simeq}
\newcommand{\vect}[1]{\overrightarrow{#1}}
\newcommand{\norme}[1]{\left\lVert #1 \right\rVert}
\newcommand{\intervalle}[2]{\left[#1\,;#2\right]}
\RequirePackage{cancel}
\renewcommand{\CancelColor}{\color{coulExo}}
\newcommand{\fois}[1]{\times{\color{coulExo}#1}}
\newcommand{\barre}[1]{\cancel{#1}}
\newlength{\lycee@eqh}\newlength{\lycee@eqd}
\newsavebox{\lycee@eqa}\newsavebox{\lycee@eqb}\newsavebox{\lycee@eqc}
\newcommand{\lycee@eqcol}[1]{\hbox{$\renewcommand{\arraystretch}{1.35}%
  \begin{array}[t]{@{}>{\displaystyle}r@{}>{\displaystyle{}}l@{}}#1\end{array}$}}
\newcommand{\lycee@eqmesure}[1]{%
  \ifdim\ht#1>\lycee@eqh\setlength{\lycee@eqh}{\ht#1}\fi
  \ifdim\dp#1>\lycee@eqd\setlength{\lycee@eqd}{\dp#1}\fi}
\newcommand{\lycee@eqtrait}{\hspace{1.2em}ou\hspace{1.2em}}
\newcommand{\lycee@eqiff}{$\iff$\hspace{0.8em}}
\newcommand{\eqproduit}[2]{\par\smallskip\noindent C'est une équation produit :
  $A\times B=0 \iff A=0$ ou $B=0$.\par\nopagebreak\smallskip
  \begingroup\setlength{\lycee@eqh}{0pt}\setlength{\lycee@eqd}{0pt}%
  \setbox\lycee@eqa=\lycee@eqcol{#1}\setbox\lycee@eqb=\lycee@eqcol{#2}%
  \lycee@eqmesure\lycee@eqa\lycee@eqmesure\lycee@eqb%
  \noindent\hspace*{\retraitcalculs}\lycee@eqiff\usebox\lycee@eqa\lycee@eqtrait\usebox\lycee@eqb\endgroup\par\smallskip}
\newcommand{\eqproduittrois}[3]{\par\smallskip\noindent C'est une équation produit :
  $A\times B\times C=0 \iff A=0$ ou $B=0$ ou $C=0$.\par\nopagebreak\smallskip
  \begingroup\setlength{\lycee@eqh}{0pt}\setlength{\lycee@eqd}{0pt}%
  \setbox\lycee@eqa=\lycee@eqcol{#1}\setbox\lycee@eqb=\lycee@eqcol{#2}\setbox\lycee@eqc=\lycee@eqcol{#3}%
  \lycee@eqmesure\lycee@eqa\lycee@eqmesure\lycee@eqb\lycee@eqmesure\lycee@eqc%
  \noindent\hspace*{\retraitcalculs}\lycee@eqiff\usebox\lycee@eqa\lycee@eqtrait\usebox\lycee@eqb\lycee@eqtrait\usebox\lycee@eqc\endgroup
  \par\smallskip}

\tikzset{grille/.style={coulGrille,line width=0.4pt}}
\newcommand{\quadrillage}[4]{\draw[grille,step=1] (#1,#2) grid (#3,#4);}
\tikzset{
  croix/.style={cross out,draw,line width=0.6pt,inner sep=0pt,minimum size=4pt},
  croixdroite/.style={croix,rotate=45,line width=0.8pt},
}
\newcommand{\pt}[3]{\path (#1) node[croix]{} node[#2,font=\small,inner sep=2.5pt]{$#3$};}

\newcommand{\lycee@repere}[4]{%
  \draw[grille,step=1] (#1,#2) grid (#3,#4);
  \draw[-{Stealth[length=2mm]},thick] (#1,0)--({#3+0.4},0) node[below left=-1pt]{$x$};
  \draw[-{Stealth[length=2mm]},thick] (0,#2)--(0,{#4+0.4}) node[below left=-1pt]{$y$};
  \node[below left,font=\scriptsize,inner sep=1.5pt] at (0,0) {$0$};}
\newcommand{\lycee@gradx}[1]{\foreach \k/\l in {#1}{%
  \draw (\k,0.07)--(\k,-0.07) node[below,font=\scriptsize,inner sep=1.5pt]{$\l$};}}
\newcommand{\lycee@grady}[1]{\foreach \k/\l in {#1}{%
  \draw (0.07,\k)--(-0.07,\k) node[left,font=\scriptsize,inner sep=1.5pt]{$\l$};}}
\newcommand{\lycee@intff}[2]{\left[#1\,;#2\right]}
\newcommand{\lycee@intfo}[2]{\left[#1\,;#2\right[}
\newcommand{\lycee@intof}[2]{\left]#1\,;#2\right]}
\newcommand{\lycee@intoo}[2]{\left]#1\,;#2\right[}
\newcommand{\lycee@ens}[1]{\left\{#1\right\}}
\AtBeginDocument{%
  \providecommand{\repere}{\lycee@repere}%
  \providecommand{\gradx}{\lycee@gradx}%
  \providecommand{\grady}{\lycee@grady}%
  \providecommand{\Cf}{\mathcal{C}\sb{f}}%
  \providecommand{\Cg}{\mathcal{C}\sb{g}}%
  \providecommand{\intff}{\lycee@intff}%
  \providecommand{\intfo}{\lycee@intfo}%
  \providecommand{\intof}{\lycee@intof}%
  \providecommand{\intoo}{\lycee@intoo}%
  \providecommand{\ens}{\lycee@ens}}

\newlist{questions}{enumerate}{3}
\setlist[questions,1]{label=\arabic*\textdegree),leftmargin=*,itemsep=2pt,topsep=3pt}
\setlist[questions,2]{label=\alph*),leftmargin=*,itemsep=1pt,topsep=2pt}
\setlist[questions,3]{label=\roman*),leftmargin=*,itemsep=1pt,topsep=1pt}
\setlist[itemize]{label=\textbullet,leftmargin=*,itemsep=2pt,topsep=3pt}

\newcommand{\besoinplace}[1]{\par\penalty\@M\begingroup
  \dimen@=#1\relax\dimen@ii=\pagegoal\advance\dimen@ii-\pagetotal
  \ifdim\dimen@>\dimen@ii\ifdim\dimen@ii>\z@\vfil\fi\break\fi\endgroup}

\newcommand{\lycee@pied}{\thepage/\pageref{LastPage}}
\newcommand{\entete}[2]{%
  \pagestyle{fancy}\fancyhf{}%
  \fancyhead[L]{\footnotesize\color{coulPart}\textbf{#1}}%
  \fancyhead[R]{\footnotesize\color{coulPart}\textbf{#2}}%
  \fancyfoot[C]{\footnotesize\lycee@pied}%
  \renewcommand{\headrulewidth}{0.4pt}%
  \renewcommand{\headrule}{\hbox to\headwidth{%
    \color{coulPart!40}\leaders\hrule height \headrulewidth\hfill}}}

\RequirePackage{ccicons}
\newcommand{\lycee@licence}{{\scriptsize\color{gray}%
  \copyright\ Luc Destombe --- \ccbyncsa\ CC BY-NC-SA 4.0}}
\newif\iflycee@licence \lycee@licencetrue
\newcommand{\sanslicence}{\lycee@licencefalse}
\AtBeginDocument{\iflycee@licence\AddToHookNext{shipout/foreground}{%
  \put(0,0){\rlap{\kern\dimexpr1in+\hoffset+\oddsidemargin+\textwidth\relax
    \llap{\raisebox{-\dimexpr1in+\voffset+\topmargin+\headheight+\headsep
      +\textheight+\footskip\relax}{\lycee@licence}}}}}\fi}

\newcommand{\lycee@site}{faire-des-maths.fr}
\newcommand{\lycee@version}{{\scriptsize\color{gray}\lycee@site
  \ifdefined\lyceeversion\ --- Version~\lyceeversion\fi}}
\AtBeginDocument{\iflycee@licence\AddToHookNext{shipout/foreground}{%
  \put(0,0){\rlap{\kern\dimexpr1in+\hoffset+\oddsidemargin\relax
    \raisebox{-\dimexpr1in+\voffset+\topmargin+\headheight+\headsep
      +\textheight+\footskip\relax}{\lycee@version}}}}\fi}

\newcommand{\formulesserrees}{%
  \setlength{\abovedisplayskip}{3pt}\setlength{\belowdisplayskip}{3pt}%
  \setlength{\abovedisplayshortskip}{2pt}\setlength{\belowdisplayshortskip}{3pt}}

\setlength{\parindent}{0pt}

\RequirePackage[implicit=false,hidelinks,pdfencoding=auto,psdextra,bookmarksopen,
  bookmarksopenlevel=1,pdfcreator={},pdfproducer={}]{hyperref}
\RequirePackage{bookmark}
\hypersetup{pdfauthor={Luc Destombe},pdfkeywords={CC BY-NC-SA 4.0}}
\newcounter{lycee@signet}
\newcommand{\signet}[3][]{\stepcounter{lycee@signet}%
  \edef\lycee@dest{\if\relax\detokenize{#1}\relax
    signet.\arabic{lycee@signet}\else#1\fi}%
  \hypertarget{\lycee@dest}{}\bookmark[level=#2,dest=\lycee@dest]{#3}}
\pdfstringdefDisableCommands{%
  \def\R{R}\def\N{N}\def\Z{Z}\def\Q{Q}\def\vect#1{#1}\def\Cf{Cf}\def\Cg{Cg}\def\fois#1{×#1}%
  \def\barre#1{#1}\def\intervalle#1#2{[#1;#2]}\def\intff#1#2{[#1;#2]}%
  \def\textsuperscript#1{#1}\def\dfrac#1#2{#1/#2}\def\frac#1#2{#1/#2}%
  \def\vec#1{#1⃗}}%
\RequirePackage[official]{eurosym}

\geometry{top=0.8cm,bottom=1.3cm,left=1cm,right=1cm,footskip=0.6cm}

\pagestyle{fancy}
\fancyhf{}
\renewcommand{\headrulewidth}{0pt}
\fancyfoot[C]{\footnotesize Page \thepage{} / \pageref{LastPage}}

\newcommand{\titrefiche}[2]{%
  \begin{center}
    \Large\textbf{#1}\\[2pt]
    \large\textbf{#2}\\[2pt]
  \end{center}
  \vspace{0.10cm}}

\setlength{\columnseprule}{0.4pt}
\setlength{\columnsep}{15pt}
\renewcommand{\columnseprulecolor}{\color{black!45}}
\setlength{\multicolsep}{2pt}

\newcounter{partiefiche}
\renewcommand{\thepartiefiche}{\Roman{partiefiche}}
\newif\ifapresbandeau
\newcommand{\lycee@bandeau}[1]{%
  \par\penalty-600
  \vspace{0.08cm}%
  \noindent\colorbox{coulPart}{%
    \parbox{\dimexpr\linewidth-2\fboxsep\relax}{%
      \color{white}\bfseries\raggedright\strut#1\strut}}%
  \nopagebreak\par\nopagebreak
  \vspace{0.06cm}\nopagebreak
  \apresbandeautrue}
\newcommand{\partiefiche}[1]{%
  \refstepcounter{partiefiche}%
  \lycee@bandeau{\signet{1}{\thepartiefiche{} --- #1}\thepartiefiche{} --- #1}}
\newcommand{\partiefichesansnum}[1]{\lycee@bandeau{\signet{1}{#1}#1}}

\newcounter{exo}
\newlength{\espaceexo}\setlength{\espaceexo}{7pt}
\newlength{\espacetitre}\setlength{\espacetitre}{2pt}
\newcommand{\exo}[1][\relax]{%
  \par
  \ifapresbandeau\apresbandeaufalse\else\penalty-150\addvspace{\espaceexo}\fi
  \refstepcounter{exo}%
  \noindent\signet[exercice-\arabic{exo}]{2}{Exercice \arabic{exo}}%
  \textbf{\color{coulExo}\underline{Exercice \theexo}}%
  \ifx\relax#1\relax\else\textbf{ : #1}\fi
  \par\nopagebreak\vspace{\espacetitre}\nopagebreak
  \ignorespaces}

\setlist[enumerate]{nosep,leftmargin=*,itemsep=2pt,topsep=2pt}
\setlist[itemize]{nosep,leftmargin=*,topsep=2pt}
\newcommand{\choix}[3]{\par\hspace*{1em}%
  \makebox[0.3\linewidth][l]{a) #1}\makebox[0.3\linewidth][l]{b) #2}c) #3}
\newcommand{\deux}[2]{\par\hspace*{0.6em}%
  \makebox[0.48\linewidth][l]{#1}#2}
\newcommand{\trois}[3]{\par\hspace*{0.6em}%
  \makebox[0.32\linewidth][l]{#1}\makebox[0.32\linewidth][l]{#2}#3}

\renewenvironment{center}{\par\addvspace{3pt}\centering}{\par\addvspace{3pt}}
\formulesserrees
\AtBeginDocument{\formulesserrees}
\clubpenalty=10000
\widowpenalty=10000
\displaywidowpenalty=10000
\brokenpenalty=10000
\setlength{\parskip}{0pt}

\RequirePackage{ragged2e}
\setlength{\RaggedRightRightskip}{0pt plus 1fil}
\AtBeginDocument{\RaggedRight\binoppenalty=10000 \relpenalty=10000 }
\g@addto@macro\@parboxrestore{\RaggedRight}
\makeatother
\usepackage{tkz-tab}
\geometry{bottom=1cm}

\setlength{\espaceexo}{4pt}
\setlist[enumerate]{nosep,leftmargin=*,itemsep=1pt,topsep=1pt}

\begin{document}

\titrefiche{Seconde --- Fonctions de référence}{Fiche d'exercices du chapitre}

{\small\itshape Le repère est orthonormé dans toute la fiche ; on trace les repères sur
le cahier, un carreau pour unité.\par}
\vspace{0.12cm}

\begin{multicols}{2}\raggedcolumns

\partiefiche{Fonction carré : courbe, variations}

\exo[Tableau de valeurs et parabole]
\begin{enumerate}
  \item Dresser le tableau de valeurs de la fonction carré de $-2$ à $2$ avec un pas
        de $0{,}5$.
  \item Tracer la parabole $y=x^{2}$ sur $\intff{-2}{2}$.
  \item Compléter par une phrase : « L'image de $-1{,}5$ par la fonction carré est
        \ldots » ; « Les antécédents de $4$ par la fonction carré sont \ldots ».
\end{enumerate}

\exo[Comparer deux carrés]
Comparer, sans les calculer :
\begin{enumerate}[label=\alph*)]
  \item $5{,}2^{2}$ et $4{,}9^{2}$
  \item $(-0{,}3)^{2}$ et $(-0{,}8)^{2}$
  \item $(-7)^{2}$ et $6{,}5^{2}$
  \item $(-2{,}4)^{2}$ et $2{,}4^{2}$
  \item $0{,}99^{2}$ et $1{,}01^{2}$
\end{enumerate}

\exo[Vrai ou faux]
Pour chaque affirmation, dire si elle est vraie ou fausse, en justifiant.
\begin{enumerate}
  \item « L'image de $-6$ par la fonction carré est $-36$. »
  \item « Si $x=4$, alors $x^{2}=16$. »
  \item « Si $x^{2}=16$, alors $x=4$. »
  \item « Les antécédents de $-9$ par la fonction carré sont $-3$ et $3$. »
  \item « Si $a<b$, alors $a^{2}<b^{2}$. »
\end{enumerate}

\partiefiche{Fonction carré : équations et inéquations}

\exo[Équations qui se ramènent à $x^{2}=k$]
Résoudre chaque équation.
\begin{enumerate}[label=\alph*)]
  \item $3x^{2}-4=x^{2}+14$
  \item $5x^{2}+3=2x^{2}+3$
  \item $7-x^{2}=3x^{2}+11$
  \item $4x^{2}+1=x^{2}+16$
  \item $2x^{2}-7=43$
  \item $(x-3)^{2}=49$
\end{enumerate}

\exo[Des aires aux longueurs]
Donner chaque longueur en valeur exacte, puis arrondie au centième.
\begin{enumerate}
  \item Une salle carrée a une aire de $98$~m$^{2}$. Quelle est la longueur de son côté ?
  \item On double l'aire d'un carré de côté $5$~cm. Quel est le côté du nouveau carré ?
        A-t-on doublé le côté ?
  \item Un disque a une aire de $49\pi$~cm$^{2}$. Quel est son rayon ?
\end{enumerate}

\exo[Chute d'un colis]
Un drone lâche un colis à $45$~m du sol. Après $t$ secondes, le colis est à la hauteur
${h(t)=45-5t^{2}}$, en mètres, pour $t$ entre $0$ et $3$.
\begin{enumerate}
  \item Calculer $h(0)$, $h(1)$ et $h(2)$.
  \item À quel instant le colis est-il à $25$~m du sol ?
  \item À quel instant touche-t-il le sol ?
  \item À quel instant est-il à $35$~m du sol ? Donner la valeur exacte, puis une
        valeur arrondie au centième.
\end{enumerate}

\exo[Lire sur la parabole]
\begin{center}
\begin{tikzpicture}[x=0.36cm,y=0.36cm]
  \repere{-4}{-1}{4}{10}
  \gradx{-3,-2,-1,1,2,3} \grady{1,2,3,4,5,6,7,8,9}
  \draw[coulCourbe,line width=1.1pt,domain=-3.1:3.1,samples=60] plot (\x,{(\x)*(\x)});
  \draw[black!70,thick] (-4,1)--(4,1) (-4,4)--(4,4) (-4,9)--(4,9);
  \node[black!70,font=\scriptsize,above] at (-3.3,1) {$y=1$};
  \node[black!70,font=\scriptsize,above] at (-3.3,4) {$y=4$};
  \node[black!70,font=\scriptsize,above] at (-3.3,9) {$y=9$};
  \node[coulCourbe,font=\footnotesize,right] at (2.9,7.4) {$y=x^{2}$};
\end{tikzpicture}
\end{center}
Résoudre graphiquement : \quad a)~$x^{2}\le 1$ \quad b)~$x^{2}>4$ \quad c)~$x^{2}=9$
\quad d)~$x^{2}\ge 9$ \quad e)~$x^{2}<0$

\exo[Résoudre exactement]
Résoudre chaque inéquation.
\begin{enumerate}[label=\alph*)]
  \item $x^{2}\le 36$
  \item $x^{2}>3$
  \item $x^{2}<-4$
  \item $x^{2}\ge 0$
  \item $4x^{2}-1\le 3x^{2}+24$
  \item $5x^{2}+2>2x^{2}+29$
\end{enumerate}

\exo[Encadrer un carré]
À l'aide des variations de la fonction carré, ou de la parabole, donner un encadrement
ou une inégalité vérifiée par $x^{2}$ dans chaque cas.

\emph{Exemple :} si $1\le x\le 3$, alors $1^{2}\le x^{2}\le 3^{2}$, soit $1\le x^{2}\le 9$,
car la fonction carré est croissante sur $\intfo{0}{+\infty}$.
\begin{enumerate}[label=\alph*)]
  \item $2\le x\le 5$
  \item $-4\le x\le -1$
  \item $-2\le x\le 3$
  \item $x>3$
  \item $x<-5$
\end{enumerate}

\exo[L'arche d'un hangar]
L'entrée d'un hangar a la forme d'une arche : la hauteur, en mètres, à l'abscisse $x$ est
$f(x)$, où $f$ est la fonction définie sur $\intff{-4}{4}$ par ${f(x)=16-x^{2}}$.
\begin{enumerate}
  \item Calculer $f(0)$ ; que représente ce nombre ?
  \item Calculer $f(-3)$ et $f(3)$. Que remarque-t-on ?
  \item Résoudre $f(x)=0$. Quelle est la largeur de l'arche au sol ?
  \item Un convoi de $7$~m de haut doit passer. Montrer que $f(x)\ge 7$ équivaut à
        $x^{2}\le 9$, puis donner la largeur du passage possible.
  \item Même question pour une hauteur d'au moins $12$~m.
\end{enumerate}

\partiefiche{Parabole et droite : position relative}

\exo[Parabole et droite]
On considère la parabole $y=x^{2}$ et la droite (d) d'équation $y=x+6$.
\begin{enumerate}
  \item Développer $(x+2)(x-3)$.
  \item Étudier le signe de $x^{2}-(x+6)$ à l'aide d'un tableau de signes.
  \item En déduire la position de la parabole et de (d).
  \item Donner les coordonnées de leurs points d'intersection.
\end{enumerate}

\exo[Un nombre et son carré]
\begin{enumerate}
  \item Pour $x\ge 0$, étudier la position de la parabole $y=x^{2}$ et de la droite
        $y=2x$.
  \item Sans calculer, comparer $0{,}7^{2}$ et $0{,}7$, puis $1{,}3^{2}$ et $1{,}3$.
  \item Léa affirme : « l'aire d'un carré, en m$^{2}$, est toujours plus grande que la
        longueur de son côté, en m ». A-t-elle raison ?
\end{enumerate}

\exo[Lire, puis démontrer]
\begin{center}
\begin{tikzpicture}[x=0.36cm,y=0.36cm]
  \repere{-4}{-2}{4}{10}
  \gradx{-3,-2,-1,1,2,3} \grady{-1,1,2,3,4,5,6,7,8,9}
  \draw[coulCourbe,line width=1.1pt,domain=-3.1:3.1,samples=60] plot (\x,{(\x)*(\x)});
  \draw[coulCourbeB,line width=1.1pt] (-4,6)--(3.5,-1.5);
  \node[coulCourbe,font=\footnotesize,right] at (2.9,7.4) {$y=x^{2}$};
  \node[coulCourbeB,font=\footnotesize,right] at (-3.9,6.9) {(d)};
\end{tikzpicture}
\end{center}
\begin{enumerate}
  \item Lire deux points de la droite (d), puis déterminer son équation réduite.
  \item Lire les coordonnées des points d'intersection de la parabole et de (d).
  \item Développer $(x+2)(x-1)$.
  \item Étudier la position de la parabole et de (d) par le calcul.
  \item En déduire les solutions de l'inéquation $x^{2}\le 2-x$.
\end{enumerate}

\partiefiche{Fonction cube}

\exo[Tableau de valeurs et courbe]
\begin{enumerate}
  \item Dresser le tableau de valeurs de la fonction cube de $-2$ à $2$ avec un pas de
        $0{,}5$.
  \item Tracer la courbe $y=x^{3}$ sur $\intff{-2}{2}$.
  \item Compléter par une phrase : « L'image de $-1{,}5$ par la fonction cube est
        \ldots » ; « L'antécédent de $-8$ par la fonction cube est \ldots ».
        Pourquoi $-8$ n'a-t-il qu'un seul antécédent ?
\end{enumerate}

\exo[Équations et inéquations]
Résoudre chaque équation ou inéquation.
\begin{enumerate}[label=\alph*)]
  \item $x^{3}=125$
  \item $x^{3}=-64$
  \item $4x^{3}-7=101$
  \item $x^{3}\ge 1$
  \item $2x^{3}+9<x^{3}+1$
  \item $x^{3}=10$ (valeur exacte, puis arrondie au centième)
\end{enumerate}

\exo[Comparer des cubes]
\begin{enumerate}
  \item Comparer, sans les calculer : \quad a)~$2{,}1^{3}$ et $1{,}9^{3}$ \quad
        b)~$(-3{,}2)^{3}$ et $(-2{,}8)^{3}$ \quad c)~$(-5)^{3}$ et $0{,}1^{3}$
  \item Comparer $(-3{,}2)^{2}$ et $(-2{,}8)^{2}$ : l'ordre est-il le même qu'au
        1.~b) ?
  \item Encadrer $x^{3}$ si $-2\le x\le 3$.
\end{enumerate}

\exo[Un aquarium cubique]
Un aquarium a la forme d'un cube d'arête $x$~dm ; son volume est $x^{3}$~dm$^{3}$, soit
$x^{3}$~litres.
\begin{enumerate}
  \item Il contient $216$~L. Quelle est son arête ?
  \item Quelle arête faut-il pour un volume d'au moins $343$~L ?
  \item On passe d'une arête de $3$~dm à une arête de $6$~dm. Par combien le volume
        est-il multiplié ?
  \item On veut un aquarium de $100$~L. Donner son arête exacte, puis arrondie au
        centième.
\end{enumerate}

\partiefiche{Fonction inverse}

\exo[Tableau de valeurs et hyperbole]
\begin{enumerate}
  \item Dresser le tableau de valeurs de la fonction inverse de $-4$ à $4$ avec un pas
        de $0{,}5$, sans $0$ ; écrire les images en fractions.
  \item Tracer l'hyperbole $y=\frac{1}{x}$ sur $\intff{-4}{4}$, sans $0$.
  \item Compléter par une phrase : « $-\frac13$ est l'image de \ldots » ; « L'antécédent
        de $\frac14$ est \ldots ». Pourquoi $0$ n'a-t-il pas d'image ?
\end{enumerate}

\exo[Équations et inéquations]
Résoudre chaque équation ; résoudre chaque inéquation à l'aide d'un tableau de signes.
\begin{enumerate}[label=\alph*)]
  \item $\dfrac{1}{x}=3$
  \item $\dfrac{1}{x}=-5$
  \item $\dfrac{1}{x}=\dfrac23$
  \item $\dfrac{1}{x}\le 4$
  \item $\dfrac{1}{x}>-2$
\end{enumerate}

\exo[Comparer des inverses]
\begin{enumerate}
  \item Comparer, sans les calculer : \quad a)~$\frac13$ et $\frac17$ \quad
        b)~$\frac{1}{-5}$ et $\frac{1}{-2}$ \quad c)~$\frac{1}{-4}$ et $\frac12$ \quad
        d)~$\frac{1}{0{,}9}$ et $\frac{1}{1{,}1}$
  \item Comparer $0{,}9^{3}$ et $1{,}1^{3}$, sans les calculer.
  \item Encadrer $\frac{1}{x}$ si $2\le x\le 5$.
\end{enumerate}

\exo[Un rectangle d'aire fixée]
Un rectangle a une aire de $12$~m$^{2}$ ; sa largeur est $x$~m ($x>0$) et sa longueur
$y$~m.
\begin{enumerate}
  \item Justifier que $y=\dfrac{12}{x}$, puis calculer $y$ pour $x=3$ et pour $x=8$.
  \item Quelle largeur donne une longueur de $8$~m ?
  \item Montrer que $y\le 6$ équivaut à $\dfrac{1}{x}\le\dfrac12$, puis en déduire les
        largeurs possibles.
  \item Calculer $(2\sqrt{3})^{2}$. En déduire que pour $x=2\sqrt{3}$, le rectangle est
        un carré.
\end{enumerate}

\partiefiche{Fonction racine carrée}

\exo[Tableau de valeurs et courbe]
\begin{enumerate}
  \item Donner les images exactes de $0$, $1$, $4$, $9$, $16$ et $25$ par la fonction
        racine carrée.
  \item Dresser le tableau de valeurs de la fonction racine carrée de $0$ à $9$ avec un
        pas de $1$, valeurs arrondies au dixième, puis tracer la courbe sur
        $\intff{0}{9}$.
  \item Compléter par une phrase : « L'antécédent de $5$ par la fonction racine carrée
        est \ldots ». Pourquoi $-4$ n'a-t-il pas d'image ?
\end{enumerate}

\exo[Équations et inéquations]
Résoudre chaque équation ou inéquation.
\begin{enumerate}[label=\alph*)]
  \item $\sqrt{x}=9$
  \item $\sqrt{x}=-4$
  \item $4\sqrt{x}-1=11$
  \item $\sqrt{x}\le 5$
  \item $\sqrt{x}>0{,}5$
  \item $\sqrt{x}<-1$
\end{enumerate}

\exo[Un nombre, son carré, sa racine]
\begin{center}
\begin{tikzpicture}[x=0.8cm,y=0.8cm]
  \repere{-0.5}{-0.5}{4}{4}
  \gradx{1,2,3} \grady{1,2,3}
  \draw[coulCourbe,line width=1.1pt,domain=0:2,samples=50] plot (\x,{(\x)*(\x)});
  \draw[coulCourbeB,line width=1.1pt] (0,0)--(3.9,3.9);
  \draw[green!50!black,line width=1.1pt,domain=0:3.9,samples=60] plot (\x,{sqrt(\x)});
  \node[coulCourbe,font=\footnotesize,left] at (1.9,3.6) {$\mathcal{C}_{1}$};
  \node[coulCourbeB,font=\footnotesize,left] at (3.4,3.6) {$\mathcal{C}_{2}$};
  \node[green!50!black,font=\footnotesize,below] at (3.4,1.75) {$\mathcal{C}_{3}$};
  \path (0,0) node[croix]{} (1,1) node[croix]{};
\end{tikzpicture}
\end{center}
Les courbes $y=x$, $y=x^{2}$ et $y=\sqrt{x}$ sont tracées ci-dessus, pour $x\ge 0$.
\begin{enumerate}
  \item Associer chaque courbe $\mathcal{C}_{1}$, $\mathcal{C}_{2}$, $\mathcal{C}_{3}$ à sa
        fonction.
  \item Ranger $x$, $x^{2}$ et $\sqrt{x}$ dans l'ordre croissant pour $0\le x\le 1$, puis
        pour $x\ge 1$.
  \item Vérifier par le calcul avec $x=0{,}25$ et avec $x=4$.
  \item Quel résultat de la partie III retrouve-t-on ?
\end{enumerate}

\exo[Vitesse d'une chute]
Un objet lâché d'une hauteur $h$, en mètres, arrive au sol à la vitesse
${v(h)=\sqrt{20h}}$, en m/s (sans frottement de l'air).
\begin{enumerate}
  \item Calculer $v(5)$ et $v(45)$.
  \item De quelle hauteur faut-il le lâcher pour qu'il arrive à $40$~m/s ?
  \item Pour quelles hauteurs la vitesse ne dépasse-t-elle pas $20$~m/s ?
  \item Sans calculer, comparer $v(10)$ et $v(11)$.
\end{enumerate}

\partiefiche{Valeur absolue d'un nombre}

\exo[Valeurs absolues et distances]
\begin{enumerate}
  \item Calculer $\lvert -9|$, $|4{,}5|$, $|0|$ et $\left|-\frac34\right|$.
  \item Calculer la distance entre : \quad a)~$2$ et $11$ \quad b)~$-5$ et $3$ \quad
        c)~$-8$ et $-2$
  \item Traduire par une phrase avec le mot « distance » : \quad a)~$|x-4|=3$ \quad
        b)~$|a+2|=5$
\end{enumerate}

\exo[Sans valeur absolue, sans racine]
Écrire chaque nombre sans valeur absolue, ni racine carrée d'un carré.
\begin{enumerate}[label=\alph*)]
  \item $|2-\sqrt{3}|$
  \item $|\sqrt{7}-3|$
  \item $|\pi-4|$
  \item $\sqrt{(1-\sqrt{2})^{2}}$
  \item $\sqrt{(-8)^{2}}$
  \item $\sqrt{(5-\sqrt{30})^{2}}$
\end{enumerate}

\exo[Priorités et valeurs absolues]
Calculer chaque expression ; une valeur absolue se calcule comme une parenthèse.
\begin{center}
$A=3-2\times\lvert -5|$ \qquad $B=|3-8|-\lvert -4|$\\[2pt]
$C=|2-7|\times\lvert -3|$ \qquad $D=\lvert -6+2|-3\times|1-3|$\\[2pt]
$E=5-|4-2\times 5|$ \qquad $F=\left|(-3)^{2}-10\right|$
\end{center}

\exo[Des écarts]
\begin{enumerate}
  \item Un matin d'hiver, il fait $-7$~°C à 6~h et $5$~°C à 14~h. Quel est l'écart de
        température ?
  \item Un plongeur est à $-18$~m (sous la surface), une mouette vole à $25$~m. Quelle
        est la distance verticale entre eux ?
  \item Sur une droite graduée, placer les nombres $x$ tels que la distance entre $x$ et
        $1$ vaut $3$. Écrire cette condition avec une valeur absolue.
\end{enumerate}

\partiefiche{Intervalle de centre $a$ et de rayon $r$}

\exo[Résoudre avec une distance]
Résoudre chaque équation ou inéquation, en s'aidant d'un schéma sur une droite graduée.
\begin{enumerate}[label=\alph*)]
  \item $|x-5|\le 1$
  \item $|x+2|\le 3$
  \item $|x-1|\le 0{,}5$
  \item $|x-4|=6$
  \item $|x+7|=2$
\end{enumerate}

\exo[Traduire un intervalle]
\emph{Exemple :} pour $\intff{2}{8}$, le centre vaut $\frac{2+8}{2}$, soit $5$, et le
rayon vaut $8-5$, soit $3$ ; donc $x\in\intff{2}{8}$ s'écrit $|x-5|\le 3$.
\begin{enumerate}
  \item Écrire chaque intervalle à l'aide d'une valeur absolue : \quad
        a)~$\intff{3}{11}$ \quad b)~$\intff{-6}{2}$ \quad c)~$\intff{-1}{0}$ \quad
        d)~$\intff{1{,}2}{1{,}8}$
  \item Résoudre $x^{2}\le 9$. Les solutions forment un intervalle : donner son centre et
        son rayon.
\end{enumerate}

\exo[Tolérances]
Écrire chaque condition avec une valeur absolue, puis donner l'intervalle
correspondant.
\begin{enumerate}
  \item Un thermostat règle la température $T$ à $19$~°C, à $1{,}5$~°C près.
  \item Une bouteille doit contenir $75$~cL, à $1$~cL près ; on note $V$ son volume.
  \item Un sondage donne $52\,\%$ des voix à un candidat, avec une marge d'erreur de
        plus ou moins $3$~points : la proportion $p$ de ses électeurs vérifie
        $|p-0{,}52|\le 0{,}03$. Est-il sûr
        d'obtenir plus de la moitié des voix ?
\end{enumerate}

\partiefiche{Fonction valeur absolue}

\exo[Courbe et lecture]
\begin{enumerate}
  \item Dresser le tableau de valeurs de la fonction valeur absolue de $-4$ à $4$ avec
        un pas de $1$, puis tracer sa courbe sur $\intff{-4}{4}$.
  \item Résoudre graphiquement : \quad a)~$|x|=2$ \quad b)~$|x|\le 2$ \quad
        c)~$|x|>3$ \quad d)~$|x|=-1$
  \item Compléter par une phrase : « Les antécédents de $2$ par la fonction valeur
        absolue sont \ldots ».
\end{enumerate}

\exo[Équations et inéquations]
Résoudre chaque équation ou inéquation.
\begin{enumerate}[label=\alph*)]
  \item $|x|=7$
  \item $|x|<4{,}5$
  \item $|x|\ge 2$
  \item $3|x|-2=10$
  \item $5-|x|>1$
  \item $|x|+3\le 1$
\end{enumerate}

\exo[Valeur absolue et carré]
\begin{center}
\begin{tikzpicture}[x=0.7cm,y=0.7cm]
  \repere{-3}{-1}{3}{4}
  \gradx{-2,-1,1,2} \grady{1,2,3}
  \draw[coulCourbe,line width=1.1pt,domain=-2:2,samples=60] plot (\x,{(\x)*(\x)});
  \draw[coulCourbeB,line width=1.1pt] (-3,3)--(0,0)--(3,3);
  \node[coulCourbe,font=\footnotesize,left] at (-1.9,3.7) {$y=x^{2}$};
  \node[coulCourbeB,font=\footnotesize,right] at (2.3,2.6) {$y=|x|$};
\end{tikzpicture}
\end{center}
\begin{enumerate}
  \item Lire les coordonnées des points d'intersection des deux courbes.
  \item Résoudre graphiquement $x^{2}\le |x|$.
  \item Pour $x\ge 0$, on a $|x|=x$ : quel résultat de la partie III démontre la réponse
        sur $\intfo{0}{+\infty}$ ?
  \item Justifier que les deux courbes sont symétriques par rapport à l'axe des
        ordonnées.
  \item Comparer $0{,}6^{2}$ et $\lvert -0{,}6|$.
\end{enumerate}

\partiefichesansnum{Problèmes de synthèse}

\exo[Le carré et le triangle]
\begin{center}
\begin{tikzpicture}[scale=0.55]
  \fill[coulCourbe!15] (0,0)--(1.5,0)--(1.5,1.5)--(0,1.5)--cycle;
  \fill[coulCourbeB!15] (1.5,0)--(4,0)--(4,4)--cycle;
  \draw[thick] (0,0)--(4,0)--(4,4)--(0,4)--cycle;
  \draw (1.5,0)--(1.5,1.5)--(0,1.5) (1.5,0)--(4,4);
  \node[below left,font=\small] at (0,0) {$A$};
  \node[below right,font=\small] at (4,0) {$B$};
  \node[above right,font=\small] at (4,4) {$C$};
  \node[above left,font=\small] at (0,4) {$D$};
  \node[below,font=\small] at (1.5,0) {$M$};
  \node[above right,font=\small] at (1.5,1.5) {$N$};
  \node[left,font=\small] at (0,1.5) {$P$};
  \draw[{Stealth[length=1.6mm]}-{Stealth[length=1.6mm]}] (0,-0.7)--(1.5,-0.7);
  \node[below,font=\small] at (0.75,-0.7) {$x$};
\end{tikzpicture}
\end{center}
$ABCD$ est un carré de côté $4$~cm. $M$ est un point de $[AB]$ et $AM=x$, avec
$0\le x\le 4$. On construit le carré $AMNP$ dans $ABCD$.
\begin{enumerate}
  \item Exprimer en fonction de $x$ l'aire $f(x)$ du carré $AMNP$, puis l'aire $g(x)$ du
        triangle $MBC$.
  \item Développer $(x+4)(x-2)$.
  \item Pour quelle position de $M$ les deux aires sont-elles égales ? Pour quelles
        valeurs de $x$ le carré a-t-il la plus grande aire ?
  \item Démontrer que $MC=\sqrt{(4-x)^{2}+16}$.
  \item Où placer $M$ pour que $MC=5$~cm ?
\end{enumerate}

\exo[Une hyperbole et une droite]
\begin{center}
\begin{tikzpicture}[x=0.6cm,y=0.6cm]
  \repere{-1}{-1}{5}{4}
  \gradx{1,2,3,4} \grady{1,2,3}
  \draw[coulCourbe,line width=1.1pt,domain=0.26:4.8,samples=60] plot (\x,{1/(\x)});
  \draw[coulCourbeB,line width=1.1pt] (0,2.5)--(3.5,-1);
  \path (0.5,2) node[croix]{} (2,0.5) node[croix]{};
  \node[coulCourbe,font=\footnotesize,above] at (4.2,0.3) {$y=\frac{1}{x}$};
  \node[coulCourbeB,font=\footnotesize,right] at (0.1,2.8) {(d)};
\end{tikzpicture}
\end{center}
On considère, pour $x>0$, l'hyperbole $y=\frac{1}{x}$ et la droite (d) d'équation
$y=-x+\frac52$.
\begin{enumerate}
  \item Vérifier que les points $E\left(\frac12\,;2\right)$ et $F\left(2\,;\frac12\right)$
        sont à la fois sur l'hyperbole et sur (d).
  \item Montrer que $\dfrac{1}{x}-\left(-x+\dfrac52\right)=\dfrac{2x^{2}-5x+2}{2x}$.
  \item Développer $(2x-1)(x-2)$.
  \item Étudier la position de l'hyperbole et de (d) sur $\intoo{0}{+\infty}$ à l'aide
        d'un tableau de signes.
  \item Quels nombres réels $x>0$ ont une somme avec leur inverse, $x+\frac{1}{x}$, qui
        ne dépasse pas $2{,}5$ ?
\end{enumerate}

\exo[Le pendule de l'horloge]
\begin{minipage}[c]{0.6\linewidth}\raggedright
Pour un pendule de longueur $L$, en mètres, on admet que la durée d'un aller-retour,
sa période, vaut ${T(L)=2\sqrt{L}}$, en secondes.
\end{minipage}\hfill
\begin{minipage}[c]{0.36\linewidth}\centering
\begin{tikzpicture}[scale=0.62]
  \foreach \k in {-1.1,-0.9,...,0.9} {\draw[black!60] (\k,0)--(\k+0.2,0.22);}
  \draw[thick] (-1.1,0)--(1.1,0);
  \draw[black!50,dashed] (0,0)--(0,-2.5);
  \draw[black!45,dashed] (0,0)--(-0.97,-2.08);
  \draw[black!45,dashed] (-0.97,-2.08) circle (0.17);
  \draw[black!60,{Stealth[length=1.6mm]}-{Stealth[length=1.6mm]}]
    (-1.1,-2.36) arc[start angle=-115,end angle=-65,radius=2.6];
  \draw[thick] (0,0)--(0.97,-2.08);
  \fill[coulCourbe] (0.97,-2.08) circle (0.17);
  \node[right,font=\small] at (0.5,-0.95) {$L$};
\end{tikzpicture}
\end{minipage}
\begin{enumerate}
  \item Calculer la période d'un pendule de $1$~m, puis de $0{,}25$~m.
  \item Quelle longueur donne une période de $3$~s ?
  \item On multiplie la longueur par $4$. Que devient la période ?
  \item Un horloger veut une période de $2$~s, à $0{,}02$~s près. Écrire cette condition
        avec une valeur absolue, puis montrer qu'elle équivaut à
        $0{,}99\le\sqrt{L}\le 1{,}01$.
  \item En déduire les longueurs possibles du pendule. À quelle précision l'horloger
        doit-il le régler ?
\end{enumerate}

\end{multicols}

\end{document}
